% Options for packages loaded elsewhere
\PassOptionsToPackage{unicode}{hyperref}
\PassOptionsToPackage{hyphens}{url}
\PassOptionsToPackage{dvipsnames,svgnames,x11names}{xcolor}
%
\documentclass[
  11pt,
]{article}
\usepackage{amsmath,amssymb}
\usepackage{iftex}
\ifPDFTeX
  \usepackage[T1]{fontenc}
  \usepackage[utf8]{inputenc}
  \usepackage{textcomp} % provide euro and other symbols
\else % if luatex or xetex
  \usepackage{unicode-math} % this also loads fontspec
  \defaultfontfeatures{Scale=MatchLowercase}
  \defaultfontfeatures[\rmfamily]{Ligatures=TeX,Scale=1}
\fi
\usepackage{lmodern}
\ifPDFTeX\else
  % xetex/luatex font selection
\fi
% Use upquote if available, for straight quotes in verbatim environments
\IfFileExists{upquote.sty}{\usepackage{upquote}}{}
\IfFileExists{microtype.sty}{% use microtype if available
  \usepackage[]{microtype}
  \UseMicrotypeSet[protrusion]{basicmath} % disable protrusion for tt fonts
}{}
\makeatletter
\@ifundefined{KOMAClassName}{% if non-KOMA class
  \IfFileExists{parskip.sty}{%
    \usepackage{parskip}
  }{% else
    \setlength{\parindent}{0pt}
    \setlength{\parskip}{6pt plus 2pt minus 1pt}}
}{% if KOMA class
  \KOMAoptions{parskip=half}}
\makeatother
\usepackage{xcolor}
\usepackage[margin=25mm]{geometry}
\usepackage{color}
\usepackage{fancyvrb}
\newcommand{\VerbBar}{|}
\newcommand{\VERB}{\Verb[commandchars=\\\{\}]}
\DefineVerbatimEnvironment{Highlighting}{Verbatim}{commandchars=\\\{\}}
% Add ',fontsize=\small' for more characters per line
\newenvironment{Shaded}{}{}
\newcommand{\AlertTok}[1]{\textcolor[rgb]{1.00,0.00,0.00}{\textbf{#1}}}
\newcommand{\AnnotationTok}[1]{\textcolor[rgb]{0.38,0.63,0.69}{\textbf{\textit{#1}}}}
\newcommand{\AttributeTok}[1]{\textcolor[rgb]{0.49,0.56,0.16}{#1}}
\newcommand{\BaseNTok}[1]{\textcolor[rgb]{0.25,0.63,0.44}{#1}}
\newcommand{\BuiltInTok}[1]{\textcolor[rgb]{0.00,0.50,0.00}{#1}}
\newcommand{\CharTok}[1]{\textcolor[rgb]{0.25,0.44,0.63}{#1}}
\newcommand{\CommentTok}[1]{\textcolor[rgb]{0.38,0.63,0.69}{\textit{#1}}}
\newcommand{\CommentVarTok}[1]{\textcolor[rgb]{0.38,0.63,0.69}{\textbf{\textit{#1}}}}
\newcommand{\ConstantTok}[1]{\textcolor[rgb]{0.53,0.00,0.00}{#1}}
\newcommand{\ControlFlowTok}[1]{\textcolor[rgb]{0.00,0.44,0.13}{\textbf{#1}}}
\newcommand{\DataTypeTok}[1]{\textcolor[rgb]{0.56,0.13,0.00}{#1}}
\newcommand{\DecValTok}[1]{\textcolor[rgb]{0.25,0.63,0.44}{#1}}
\newcommand{\DocumentationTok}[1]{\textcolor[rgb]{0.73,0.13,0.13}{\textit{#1}}}
\newcommand{\ErrorTok}[1]{\textcolor[rgb]{1.00,0.00,0.00}{\textbf{#1}}}
\newcommand{\ExtensionTok}[1]{#1}
\newcommand{\FloatTok}[1]{\textcolor[rgb]{0.25,0.63,0.44}{#1}}
\newcommand{\FunctionTok}[1]{\textcolor[rgb]{0.02,0.16,0.49}{#1}}
\newcommand{\ImportTok}[1]{\textcolor[rgb]{0.00,0.50,0.00}{\textbf{#1}}}
\newcommand{\InformationTok}[1]{\textcolor[rgb]{0.38,0.63,0.69}{\textbf{\textit{#1}}}}
\newcommand{\KeywordTok}[1]{\textcolor[rgb]{0.00,0.44,0.13}{\textbf{#1}}}
\newcommand{\NormalTok}[1]{#1}
\newcommand{\OperatorTok}[1]{\textcolor[rgb]{0.40,0.40,0.40}{#1}}
\newcommand{\OtherTok}[1]{\textcolor[rgb]{0.00,0.44,0.13}{#1}}
\newcommand{\PreprocessorTok}[1]{\textcolor[rgb]{0.74,0.48,0.00}{#1}}
\newcommand{\RegionMarkerTok}[1]{#1}
\newcommand{\SpecialCharTok}[1]{\textcolor[rgb]{0.25,0.44,0.63}{#1}}
\newcommand{\SpecialStringTok}[1]{\textcolor[rgb]{0.73,0.40,0.53}{#1}}
\newcommand{\StringTok}[1]{\textcolor[rgb]{0.25,0.44,0.63}{#1}}
\newcommand{\VariableTok}[1]{\textcolor[rgb]{0.10,0.09,0.49}{#1}}
\newcommand{\VerbatimStringTok}[1]{\textcolor[rgb]{0.25,0.44,0.63}{#1}}
\newcommand{\WarningTok}[1]{\textcolor[rgb]{0.38,0.63,0.69}{\textbf{\textit{#1}}}}
\usepackage{longtable,booktabs,array}
\usepackage{calc} % for calculating minipage widths
% Correct order of tables after \paragraph or \subparagraph
\usepackage{etoolbox}
\makeatletter
\patchcmd\longtable{\par}{\if@noskipsec\mbox{}\fi\par}{}{}
\makeatother
% Allow footnotes in longtable head/foot
\IfFileExists{footnotehyper.sty}{\usepackage{footnotehyper}}{\usepackage{footnote}}
\makesavenoteenv{longtable}
\setlength{\emergencystretch}{3em} % prevent overfull lines
\providecommand{\tightlist}{%
  \setlength{\itemsep}{0pt}\setlength{\parskip}{0pt}}
\setcounter{secnumdepth}{5}
\ifLuaTeX
  \usepackage{selnolig}  % disable illegal ligatures
\fi
\usepackage{bookmark}
\IfFileExists{xurl.sty}{\usepackage{xurl}}{} % add URL line breaks if available
\urlstyle{same}
\hypersetup{
  pdftitle={Integral Applications: Study Guide},
  colorlinks=true,
  linkcolor={blue},
  filecolor={Maroon},
  citecolor={Blue},
  urlcolor={blue},
  pdfcreator={LaTeX via pandoc}}

\title{Integral Applications: Study Guide}
\author{}
\date{15 October 2026}

\begin{document}
\maketitle

{
\hypersetup{linkcolor=}
\setcounter{tocdepth}{3}
\tableofcontents
}
\section{Concrete Lebesgue and Hausdorff integral
applications}\label{start}

A first collection of \textbf{16 fully worked examples}, written for an
analysis student. Start with the \hyperref[problems]{problem index},
then open the separate solution that you want. The examples favor
explicit computations, stated hypotheses, and direct arguments
explaining why the integration theorem applies.

\begin{itemize}
\tightlist
\item
  \hyperref[conventions]{Measures and theorem conventions}: normalized
  Hausdorff measure, completed-product sections, signed integrability,
  area and coarea.
\item
  \hyperref[problems]{Problem index}: no computed answers.
\item
  \hyperref[atlas]{Counterexample atlas}: precise failures and
  sufficient conditions for exchanging sums, limits, derivatives, and
  integrals.
\item
  \href{https://github.com/Drsunflowerherbert/Codex-for-Math/tree/codex/integral-examples-2026-10-15/integral-examples/pdf/integral-examples.pdf}{Full
  worked collection, PDF} and
  \href{https://github.com/Drsunflowerherbert/Codex-for-Math/tree/codex/integral-examples-2026-10-15/integral-examples/pdf/integral-examples.tex}{editable
  LaTeX}.
\item
  \href{https://github.com/Drsunflowerherbert/Codex-for-Math/tree/codex/integral-examples-2026-10-15/integral-examples/pdf/guide.pdf}{Compact
  guide, PDF} and
  \href{https://github.com/Drsunflowerherbert/Codex-for-Math/tree/codex/integral-examples-2026-10-15/integral-examples/pdf/guide.tex}{LaTeX}.
\item
  \hyperref[bibliography]{Source references},
  \href{https://github.com/Drsunflowerherbert/Codex-for-Math/tree/codex/integral-examples-2026-10-15/integral-examples/checks/source-references.json}{source
  snapshot hashes}, and
  \href{https://github.com/Drsunflowerherbert/Codex-for-Math/tree/codex/integral-examples-2026-10-15/integral-examples/VERIFICATION.md}{verification
  ledger}.
\item
  \href{https://github.com/Drsunflowerherbert/Codex-for-Math/tree/codex/integral-examples-2026-10-15/integral-examples/index.json}{Machine-readable
  index} without results, and
  \href{https://github.com/Drsunflowerherbert/Codex-for-Math/tree/codex/integral-examples-2026-10-15/integral-examples/metadata.json}{full
  metadata} with hypotheses, results, and review status.
\end{itemize}

\subsection{Find a technique again}\label{find-a-technique-again}

From the repository root, using Python 3 with its standard library:

\begin{Shaded}
\begin{Highlighting}[]
\ExtensionTok{python3}\NormalTok{ integral{-}examples/find\_examples.py }\AttributeTok{{-}{-}topic}\NormalTok{ coarea}
\ExtensionTok{python3}\NormalTok{ integral{-}examples/find\_examples.py }\AttributeTok{{-}{-}technique}\NormalTok{ multiplicity}
\ExtensionTok{python3}\NormalTok{ integral{-}examples/find\_examples.py }\AttributeTok{{-}{-}prerequisite}\NormalTok{ dominated}
\ExtensionTok{python3}\NormalTok{ integral{-}examples/find\_examples.py }\AttributeTok{{-}{-}text}\NormalTok{ concentration}
\ExtensionTok{python3}\NormalTok{ integral{-}examples/find\_examples.py }\AttributeTok{{-}{-}list{-}tags}
\end{Highlighting}
\end{Shaded}

Filters are case-insensitive substring matches and combine with AND.
They print titles, tags, and paths relative to
\texttt{integral-\allowbreak{}examples/\allowbreak{}}. The \texttt{-\/-text} filter searches IDs,
titles, paths, topics, prerequisites, and techniques; it does not search
solution prose or recorded hypotheses/results. Add
\texttt{-\/-solutions} to print the recorded result. The files remain
ordinary Markdown and JSON; retrieval needs no hosted service.

\subsection{How to study with the
collection}\label{how-to-study-with-the-collection}

For each problem, identify the function, its domain, the measure, and
the integrability question before computing. In a region problem, write
the indicator and determine its sections. For signed integrals, examine
the absolute value or the positive and negative parts. For a
parametrization, compute the Jacobian and count preimages before
identifying the answer with an integral over the image. For a parameter
limit, write the proposed majorant explicitly.

E01--E06 cover Lebesgue integration and iterated integrals; E07--E13
cover curves, surfaces, area, coarea, and singular maps; E14--E16 cover
limit and derivative exchanges and approximate identities. Cross-links
in the atlas make failure cases easy to retrieve alongside their
positive criteria.

\subsection{Reproduce the checks and
PDFs}\label{reproduce-the-checks-and-pdfs}

\begin{Shaded}
\begin{Highlighting}[]
\ExtensionTok{python3}\NormalTok{ integral{-}examples/checks/check\_examples.py}
\ExtensionTok{python3}\NormalTok{ integral{-}examples/checks/check\_collection.py}
\ExtensionTok{python3}\NormalTok{ integral{-}examples/build.py}
\end{Highlighting}
\end{Shaded}

The concrete identity checker uses SymPy and mpmath. The collection
checker and retrieval command use only the Python standard library. The
build uses Pandoc and pdfLaTeX. See
\href{https://github.com/Drsunflowerherbert/Codex-for-Math/tree/codex/integral-examples-2026-10-15/integral-examples/checks/requirements.txt}{checks/requirements.txt}
for the versions actually used. The checked-in PDF and LaTeX need no
build tools to read.

Independent model reviewers audited every example; the ledger records
the scope, findings, and revisions. Symbolic and numerical checks
supplement those reviews and are \textbf{not proofs of analytic
hypotheses}. No Lean verification or human peer review is claimed.
General measure-theoretic theorems are cited prerequisites; this
collection does not reprove all of measure theory.

All collection changes are confined to \texttt{integral-\allowbreak{}examples/\allowbreak{}} on
branch \texttt{codex/\allowbreak{}integral-\allowbreak{}examples-\allowbreak{}2026-\allowbreak{}10-\allowbreak{}15}, based on the
integrated research commit
\texttt{e4f337c3\allowbreak{}aa52ec4c\allowbreak{}3f2aab07\allowbreak{}8d7c91b3\allowbreak{}2ee0015f}. Other research
branches are independent. The earlier paper was not used or
reconstructed. No third-party problem collection or source PDF is
reproduced, and no license is assigned here to the cited sources or to
unrelated repository work.

\section{Problem index --- without computed answers}\label{problems}

Each entry links to a separate full solution. Read
\hyperref[conventions]{CONVENTIONS.md} first for measure normalization
and the exact role of almost-everywhere sections. The JSON index and
retrieval command described in \hyperref[start]{README.md} support
search by topic, prerequisite, and technique without printing answers by
default.

\begin{longtable}[]{@{}
  >{\raggedright\arraybackslash}p{(\columnwidth - 6\tabcolsep) * \real{0.0700}}
  >{\raggedright\arraybackslash}p{(\columnwidth - 6\tabcolsep) * \real{0.3700}}
  >{\raggedright\arraybackslash}p{(\columnwidth - 6\tabcolsep) * \real{0.2500}}
  >{\raggedright\arraybackslash}p{(\columnwidth - 6\tabcolsep) * \real{0.3100}}@{}}
\toprule\noalign{}
\begin{minipage}[b]{\linewidth}\raggedright
ID
\end{minipage} & \begin{minipage}[b]{\linewidth}\raggedright
Task
\end{minipage} & \begin{minipage}[b]{\linewidth}\raggedright
Main prerequisites
\end{minipage} & \begin{minipage}[b]{\linewidth}\raggedright
Technique to look for
\end{minipage} \\
\midrule\noalign{}
\endhead
\bottomrule\noalign{}
\endlastfoot
\href{https://github.com/Drsunflowerherbert/Codex-for-Math/tree/codex/integral-examples-2026-10-15/integral-examples/examples/E01-triangular-region.md}{E01}
& Integrate \(e^{y^2}\) over \(0<x<y<1\). & Measurability, Tonelli &
Change order on a measurable region \\
\href{https://github.com/Drsunflowerherbert/Codex-for-Math/tree/codex/integral-examples-2026-10-15/integral-examples/examples/E02-positive-measure-boundary.md}{E02}
& Construct a region whose boundary has positive area; integrate \(2y\)
over it. & Countable additivity, Tonelli & Explicit removed-interval
sums \\
\href{https://github.com/Drsunflowerherbert/Codex-for-Math/tree/codex/integral-examples-2026-10-15/integral-examples/examples/E03-exceptional-sections.md}{E03}
& Can every section of a Lebesgue-measurable null set be integrated?
Prove the correct general statement. & Completion, Fubini & Borel
representative and null envelope \\
\href{https://github.com/Drsunflowerherbert/Codex-for-Math/tree/codex/integral-examples-2026-10-15/integral-examples/examples/E04-signed-fubini-trap.md}{E04}
& Compare the iterated integrals of \((x^2-y^2)/(x^2+y^2)^2\) on
\((0,1)^2\). & Signed integrals, Fubini & Check positive and negative
parts first \\
\href{https://github.com/Drsunflowerherbert/Codex-for-Math/tree/codex/integral-examples-2026-10-15/integral-examples/examples/E05-signed-fubini-counting.md}{E05}
& Sum an explicit signed matrix by rows and by columns. & Counting
measure, series & Detect infinite absolute mass \\
\href{https://github.com/Drsunflowerherbert/Codex-for-Math/tree/codex/integral-examples-2026-10-15/integral-examples/examples/E06-frullani-parameter-integral.md}{E06}
& Evaluate and differentiate a Frullani parameter integral. & Tonelli,
dominated convergence & Replace a quotient by a parameter integral \\
\href{https://github.com/Drsunflowerherbert/Codex-for-Math/tree/codex/integral-examples-2026-10-15/integral-examples/examples/E07-parabola.md}{E07}
& Integrate the first coordinate along a parabola with respect to
length. & Derivatives, area formula & Compute the speed and pull back
the weight \\
\href{https://github.com/Drsunflowerherbert/Codex-for-Math/tree/codex/integral-examples-2026-10-15/integral-examples/examples/E08-circle-multiplicity.md}{E08}
& Integrate a coordinate square on a twice-traversed circle. & Curve
area formula & Count multiplicity; distinguish circulation \\
\href{https://github.com/Drsunflowerherbert/Codex-for-Math/tree/codex/integral-examples-2026-10-15/integral-examples/examples/E09-paraboloid-graph.md}{E09}
& Integrate height on a paraboloid graph over the unit disk. & Gram
determinant, polar coordinates & Derive the graph area factor \\
\href{https://github.com/Drsunflowerherbert/Codex-for-Math/tree/codex/integral-examples-2026-10-15/integral-examples/examples/E10-sphere-moment.md}{E10}
& Integrate a coordinate square over a sphere of radius \(R\). & Surface
area formula & Handle chart seams and pole exceptions \\
\href{https://github.com/Drsunflowerherbert/Codex-for-Math/tree/codex/integral-examples-2026-10-15/integral-examples/examples/E11-fold-area-multiplicity.md}{E11}
& Apply weighted area to a folded square parametrization. & Jacobian,
area formula & Compare source integral and image multiplicity \\
\href{https://github.com/Drsunflowerherbert/Codex-for-Math/tree/codex/integral-examples-2026-10-15/integral-examples/examples/E12-radial-coarea-critical-value.md}{E12}
& Turn an integral on a disk into integrals over circular level sets. &
Gradient, coarea & Account for a critical point before division \\
\href{https://github.com/Drsunflowerherbert/Codex-for-Math/tree/codex/integral-examples-2026-10-15/integral-examples/examples/E13-cantor-singular-map.md}{E13}
& Explain why the Cantor staircase violates a proposed derivative-area
identity. & Cantor set, absolute continuity & Locate the failed
Lipschitz hypothesis \\
\href{https://github.com/Drsunflowerherbert/Codex-for-Math/tree/codex/integral-examples-2026-10-15/integral-examples/examples/E14-concentration-escape-atlas.md}{E14}
& Track mass in concentration and escape examples when taking limits
under integrals. & Pointwise convergence, DCT & Uniform integrability
and spatial tightness \\
\href{https://github.com/Drsunflowerherbert/Codex-for-Math/tree/codex/integral-examples-2026-10-15/integral-examples/examples/E15-differentiation-concentration.md}{E15}
& Compare a parameter derivative with the integral of pointwise
derivatives. & Difference quotients, DCT & Test local domination;
distinguish derivative limits \\
\href{https://github.com/Drsunflowerherbert/Codex-for-Math/tree/codex/integral-examples-2026-10-15/integral-examples/examples/E16-exponential-approximate-identity.md}{E16}
& Recover a boundary value from exponential averages. & Substitution,
DCT & Approximate identity and an explicit tail estimate \\
\end{longtable}

Suggested routes: E01--E06 for Lebesgue integration and iterated
integrals; E07--E13 for geometric integration; E14--E16 for exchanges of
limits, derivatives, and integrals. E03 and E04 are useful safeguards
before generalizing a successful calculation. The
\hyperref[atlas]{counterexample atlas} maps the safeguards to precise
sufficient conditions and contains some outcome spoilers.

\section{Measures, functions, and theorems used in this
library}\label{conventions}

The examples are applications for an analysis student. Their proofs are
written out; the general integration theorems below are prerequisites,
with sources in \hyperref[bibliography]{SOURCES.md}. A computation check
verifies an identity, not the measurability or integrability hypotheses
that justify it.

\subsection{Functions and values}\label{functions-and-values}

A function is a map with a specified domain, codomain, and rule. For
example, \(f:(0,1)^2\to\mathbb R\) is a function, \(f(x,y)\) is its
value at \((x,y)\), and \(f_x:y\mapsto f(x,y)\) is a section function.
With a fixed integrand and measure, the integral is a number (or an
allowed extended value), not another notation for the function. A
parameter integral defines a new function, for example
\(F:t\mapsto\int f_t\,d\mu\); \(F(t)\) is its value. When a section
integral is defined only almost everywhere, a measurable version is
chosen before integrating that version in the outer variable.

\subsection{Lebesgue and counting
measures}\label{lebesgue-and-counting-measures}

\(\mathcal L^n\) denotes complete \(n\)-dimensional Lebesgue measure on
\(\mathbb R^n\). A restricted integral \(\int_E f\,d\mathcal L^n\) means
integration over the measurable set \(E\) with the restricted measure.
For finite-valued \(f\) defined on all of \(\mathbb R^n\), it equals
\(\int f\mathbf1_E\,d\mathcal L^n\). Counting measure on
\(\mathbb N=\{1,2,\ldots\}\) turns integration into an explicit sum. The
Lebesgue and countable counting measure factors used in the product
theorems here are sigma-finite. This does not assert sigma-finiteness of
\(\mathcal H^k\) on all of \(\mathbb R^n\) when \(k<n\); our curve and
surface examples use finite-measure or countably covered restrictions.

For real measurable \(f\), set \(f^+=\max(f,0)\) and \(f^-=\max(-f,0)\).
A finite Lebesgue integral requires

\[
\int |f|\,d\mu=\int f^+\,d\mu+\int f^-\,d\mu<\infty.
\]

An extended signed integral is defined if at least one of the two terms
is finite. If both are infinite, their difference is undefined. An
improper limit or principal value must be named separately; its
existence does not supply a signed Lebesgue integral.

\subsection{Normalized Hausdorff
measure}\label{normalized-hausdorff-measure}

For an integer \(k\geq1\), let \(\omega_k=\mathcal L^k(B(0,1))\). We use

\[
\mathcal H^k(E)=\lim_{\delta\downarrow0}\inf\left\{
\sum_{i=1}^{\infty}\omega_k\left(\frac{\operatorname{diam}U_i}{2}\right)^k:
E\subseteq\bigcup_{i=1}^{\infty}U_i,\quad
\operatorname{diam}U_i\leq\delta\right\}.
\]

\(\mathcal H^0\) is counting measure. This normalization makes
\(\mathcal H^k\) agree with \(\mathcal L^k\) on every \(k\)-plane, using
an isometry to identify that plane with \(\mathbb R^k\). Thus
\(\mathcal H^1\) measures length and \(\mathcal H^2\) measures surface
area. In particular \(\omega_1=2\) and \(\omega_2=\pi\). Some books omit
the factor \(\omega_k/2^k\); their numerical values must be rescaled
before using these examples. All Hausdorff integrals in this library use
the convention in this paragraph.

\subsection{Tonelli and Fubini, including
completion}\label{tonelli-and-fubini-including-completion}

For a nonnegative product-measurable function on sigma-finite measure
spaces, Tonelli permits either order of integration, including the value
\(+\infty\). For a signed or complex product-measurable function, Fubini
applies if the absolute value has finite integral. It then gives finite
section integrals almost everywhere and equality of the two iterated
integrals with the product integral.

The product of two complete Lebesgue sigma-algebras need not be
complete. For \(\mathcal L^{p+q}\)-measurable sets or functions on
\(\mathbb R^{p+q}\), use the completed-product versions: sections need
be measurable only for almost every parameter. A Borel set has Borel
sections at every parameter; an arbitrary Lebesgue-measurable set need
not have measurable sections at every parameter. The outer integral uses
a measurable representative of the almost-everywhere section-integral
function. E03 proves this distinction directly with a null envelope. No
claim about all sections is inferred from completeness.

\subsection{Area with multiplicity}\label{area-with-multiplicity}

For a Lipschitz map \(F:\mathbb R^m\to\mathbb R^n\) with \(m\leq n\),
write

\[
J_mF(x)=\sqrt{\det(DF(x)^TDF(x))},\qquad
N(F,A,z)=\#\{x\in A:F(x)=z\}.
\]

The derivative and Jacobian exist almost everywhere; set the Jacobian to
zero at the exceptional points where the derivative does not exist. For
a measurable domain \(A\) and nonnegative Borel weight \(h\) on the
source, the weighted area formula is

\[
\int_A h(x)J_mF(x)\,d\mathcal L^m(x)
=\int_{\mathbb R^n}\left(\sum_{x\in A\cap F^{-1}(\{z\})}h(x)\right)
\,d\mathcal H^m(z).
\]

The sum is an extended nonnegative sum. With \(h=1\) it counts
multiplicity. With a nonnegative Borel \(w\) and \(h=w\circ F\) it
becomes \(w(z)N(F,A,z)\), with the nonnegative-integration convention
\(0\cdot\infty=0\). An injective parametrization therefore computes the
geometric scalar integral once. If multiplicity is \(q\) almost
everywhere, the parameter integral counts it \(q\) times. For signed
weights require finite absolute integral with this multiplicity. A map
on a bounded domain that is Lipschitz there can be extended to apply the
formula; the geometric parametrizations to which we apply area are also
Lipschitz on a suitable containing compact box. The Cantor staircase is
an explicitly excluded counterexample.

For a graph \(F(u)=(u,g(u))\) with a scalar \(C^1\) function \(g\),
\(J_mF=\sqrt{1+|\nabla g|^2}\). This factor is derived in the graph
example.

\subsection{Coarea and division by a
Jacobian}\label{coarea-and-division-by-a-jacobian}

For a Lipschitz map \(G:\mathbb R^n\to\mathbb R^m\), \(n\geq m\), the
coarea Jacobian is

\[
J_mG(x)=\sqrt{\det(DG(x)DG(x)^T)}.
\]

For nonnegative Borel \(h\), coarea says

\[
\int_{\mathbb R^n}h(x)J_mG(x)\,d\mathcal L^n(x)
=\int_{\mathbb R^m}\left(\int_{G^{-1}(\{t\})}h(x)
\,d\mathcal H^{n-m}(x)\right)\,d\mathcal L^m(t).
\]

Fiber statements and inner integrals have their usual almost-everywhere
meaning. Here \(m=1\) gives \(J_1G=|\nabla G|\). To recover an
unweighted volume integral, division by \(J_mG\) is legitimate only on
its positive set. Any contribution from its zero set must be handled
separately. E12 identifies and disposes of its critical point
explicitly. E13 shows why a singular continuous map cannot be
substituted into a Lipschitz area formula.

\subsection{Scalar integration and
orientation}\label{scalar-integration-and-orientation}

\(\int_S w\,d\mathcal H^k\) has no orientation. A scalar curve
parametrization uses speed \(|\gamma'(t)|\). A vector circulation uses
\(\int F(\gamma(t))\cdot\gamma'(t)\,dt\); reversing traversal changes
its sign. Flux through an oriented surface is
\(\int_S F\cdot\nu\,d\mathcal H^{n-1}\) with a chosen unit normal
\(\nu\). Changing \(\nu\) changes flux. Neither a tangent vector nor a
normal vector is part of Hausdorff measure itself.

\subsection{Convergence and
differentiation}\label{convergence-and-differentiation}

The version of monotone convergence used here concerns nonnegative
increasing measurable functions. Dominated convergence concerns
measurable functions and requires almost-everywhere convergence and one
integrable majorant for the entire sequence. Pointwise convergence,
boundedness of each individual function, or a uniform bound on the
integrals alone is insufficient. E14 gives concentration and escape
examples with explicit sufficient repairs. E15 states and proves a local
domination criterion for differentiating an integral. E16 explains how
concentration can instead be used deliberately as an approximate
identity.

\section{A small atlas of failed exchanges}\label{atlas}

This is a retrieval map to complete proofs, not a replacement for them.
It gives outcomes; use \hyperref[problems]{PROBLEMS.md} for the
problem-first route.

\begin{longtable}[]{@{}
  >{\raggedright\arraybackslash}p{(\columnwidth - 4\tabcolsep) * \real{0.2800}}
  >{\raggedright\arraybackslash}p{(\columnwidth - 4\tabcolsep) * \real{0.3600}}
  >{\raggedright\arraybackslash}p{(\columnwidth - 4\tabcolsep) * \real{0.3600}}@{}}
\toprule\noalign{}
\begin{minipage}[b]{\linewidth}\raggedright
Exchange or tempting inference
\end{minipage} & \begin{minipage}[b]{\linewidth}\raggedright
Example and failure mechanism
\end{minipage} & \begin{minipage}[b]{\linewidth}\raggedright
Positive sufficient condition
\end{minipage} \\
\midrule\noalign{}
\endhead
\bottomrule\noalign{}
\endlastfoot
Integrate every section of a Lebesgue-measurable set & E03: a
nonmeasurable set can sit in one exceptional section of a planar null
set. & Borel measurability gives Borel sections everywhere; completed
Lebesgue measurability gives measurable sections almost everywhere. \\
Swap signed iterated integrals & E04: both signed parts have infinite
integral near the origin. E05: the signed matrix has infinite absolute
sum. & Joint measurability and finite integral of the absolute value
give Fubini; nonnegativity gives Tonelli with extended values. \\
Pass an almost-everywhere limit through an integral & E14: mass
concentrates at a shrinking interval; its height grows. & DCT with one
integrable majorant; on a finite measure space, convergence in measure
plus uniform integrability gives convergence in \(L^1\). \\
Treat bounded integral mass as sufficient on an infinite space & E14: a
unit bump moves to infinity. Value-tail uniform integrability alone
permits this. & Add spatial tightness to value-tail uniform
integrability, together with convergence in measure on bounded sets (and
an integrable limit). Precise versions are stated in E14. \\
Differentiate a parameter integral at an endpoint & E15:
difference-quotient mass concentrates near zero; pointwise derivatives
miss that mass. & Measurability, pointwise differentiability almost
everywhere, one integrable local bound on the parameter derivative, and
one integrable parameter slice. See the exact interval and endpoint
conditions in E15. \\
Differentiate a uniform limit of differentiable functions & E15 related
example: \(\sin(nx)/n\) tends uniformly to zero while its derivative is
\(1\) at zero for every \(n\). & On a compact interval, uniform
convergence of the derivatives and convergence of the function values at
one point imply a differentiable limit with the expected derivative. The
mean value theorem proves this for differentiable functions; for \(C^1\)
functions one can also use the fundamental theorem of calculus. \\
Replace geometric area by parameter area & E08/E11: repeated or folded
parametrizations count the image more than once. & Injectivity outside
sets whose contributions are zero, or explicit multiplicity in the area
formula. \\
Divide by a coarea Jacobian everywhere & E12: the gradient vanishes at
the origin. & Restrict division to the positive-Jacobian set, then prove
the zero set contributes zero or account for its contribution
separately. \\
Use the area formula for a merely continuous monotone map & E13: the
Cantor staircase has derivative zero almost everywhere but maps onto an
interval. & Lipschitz maps satisfy the area formula; in one dimension
absolute continuity gives the fundamental theorem and the relevant
change-of-variable identities. \\
\end{longtable}

Concentration is also a tool: E16 uses an approximate identity. There
the integration variable is rescaled and a fixed integrable majorant
becomes available. The conclusion is a boundary evaluation, with all its
hypotheses stated, rather than an exchange based solely on pointwise
convergence in the original variable.

\section{Sources and conventions audit}\label{bibliography}

All links below were opened and checked on \textbf{2026-10-15}. These
are author-hosted teaching texts, used for theorem statements and
conventions. The worked examples in this directory are original
expositions; the linked texts' proofs and problem collections are not
reproduced. PDF page numbers below count the cover as page 1. A PDF
viewer may display a different logical page label. Online files can
change.

\subsection{\texorpdfstring{\texttt{hunter-\allowbreak{}measure}}{hunter-measure}}\label{hunter-measure}

John K. Hunter, \emph{Measure Theory}, University of California, Davis,
2011.
\href{https://www.math.ucdavis.edu/~hunter/measure_theory/measure_notes.pdf}{Author's
PDF}.

\begin{itemize}
\tightlist
\item
  Monotone convergence: Theorem 4.6, printed p.41, PDF p.45.
\item
  Dominated convergence: Theorem 4.24, printed p.49,
  \href{https://www.math.ucdavis.edu/~hunter/measure_theory/measure_notes.pdf\#page=53}{PDF
  p.53}; a.e. versions are discussed in §4.6, printed p.50.
\item
  Product-measurable sections: §5.1, Proposition 5.2, printed p.55, PDF
  p.59; function sections: Theorem 5.15, printed p.60.
\item
  Fubini with absolute integrability: Theorem 5.18, printed p.61,
  \href{https://www.math.ucdavis.edu/~hunter/measure_theory/measure_notes.pdf\#page=65}{PDF
  p.65}.
\item
  Completion and exceptional sections: Example 5.20 and Theorem 5.21,
  printed p.62,
  \href{https://www.math.ucdavis.edu/~hunter/measure_theory/measure_notes.pdf\#page=66}{PDF
  p.66}.
\end{itemize}

The completed Lebesgue product statement guarantees measurable and
integrable sections almost everywhere. It does not guarantee
measurability of every section. Theorem 5.15's sentence defining complex
section integrals suppresses existence hypotheses; use the explicit
integrability hypotheses of 5.18/5.21.

\subsection{\texorpdfstring{\texttt{axler-\allowbreak{}mira}}{axler-mira}}\label{axler-mira}

Sheldon Axler, \emph{Measure, Integration \& Real Analysis}, Graduate
Texts in Mathematics 282, Springer, 2020; author's online revision dated
12 June 2026. \href{https://measure.axler.net/MIRA.pdf}{Author's PDF}.

\begin{itemize}
\tightlist
\item
  Tonelli: §5B, Theorem 5.28, printed p.129,
  \href{https://measure.axler.net/MIRA.pdf\#page=144}{PDF p.144}.
\item
  Nonnegative double sums: Corollary 5.31, printed p.131, PDF p.146.
\item
  Fubini: Theorem 5.32, printed p.132,
  \href{https://measure.axler.net/MIRA.pdf\#page=147}{PDF p.147}, with
  exceptional-parameter conventions immediately above and in the proof.
\end{itemize}

These product theorems assume sigma-finite factors. Tonelli accepts
nonnegative product-measurable functions and permits the value infinity;
Fubini requires a finite integral of the absolute value. This is an
independent reference for the distinction used in the examples. For
completed Lebesgue measurability rather than product measurability, also
use \texttt{hunter-\allowbreak{}measure} §5.7.

\subsection{\texorpdfstring{\texttt{simon-\allowbreak{}gmt}}{simon-gmt}}\label{simon-gmt}

Leon Simon, \emph{Introduction to Geometric Measure Theory}, Stanford
University, 2017/2018 NTU lecture draft; last-update line: 7 March 2018.
\href{https://math.stanford.edu/~lms/ntu-gmt-text.pdf}{Author's PDF}.

\begin{itemize}
\tightlist
\item
  Normalization: Chapter 1, definitions 2.1--2.2, printed p.11,
  \href{https://math.stanford.edu/~lms/ntu-gmt-text.pdf\#page=9}{PDF
  p.9}; equality with Lebesgue outer measure: Theorem 2.6, printed p.12,
  PDF p.10.
\item
  Area and multiplicity: Chapter 2, Theorem 3.3, printed p.57,
  \href{https://math.stanford.edu/~lms/ntu-gmt-text.pdf\#page=32}{PDF
  p.32}; weighted fiber sum: Remark 3.4(2), printed p.58, PDF p.33;
  curves and submanifolds: Examples 3.5, printed pp.60--61, PDF
  pp.34--35.
\item
  Coarea: Chapter 2, Theorem 7.3 and weighted Remark 7.4(1), printed
  p.76,
  \href{https://math.stanford.edu/~lms/ntu-gmt-text.pdf\#page=42}{PDF
  p.42}.
\item
  Critical parts of fibers: Theorem 7.5, printed p.80, PDF p.44; the
  stronger smooth Sard conclusion is distinguished in Remark 7.6.
\end{itemize}

The PDF has two printed pages per sheet. This version uses the factor
\texttt{omega\_k\ /\ 2\^{}k} multiplying sums of diameter powers, where
\texttt{omega\_k} is unit-ball volume. Consequently normalized Hausdorff
measure agrees with Lebesgue measure on k-planes. Counting measure is
the k=0 convention.

\subsection{\texorpdfstring{\texttt{hajlasz-\allowbreak{}geometry}}{hajlasz-geometry}}\label{hajlasz-geometry}

Piotr Hajłasz, \emph{Geometric Analysis}, University of Pittsburgh,
undated author-hosted lecture notes.
\href{https://www.pitt.edu/~hajlasz/Notatki/Analysis_4.pdf}{Author's
PDF}.

\begin{itemize}
\tightlist
\item
  Normalization: §3.1, printed/PDF p.19; equality with Lebesgue measure:
  Theorem 3.2, p.20.
\item
  Surface integration: Theorem 3.17 and following discussion, p.25.
\item
  Area and multiplicity with the Lusin N condition: Theorem 3.18 and
  Remark 3.19, pp.26--27.
\item
  Coarea: §3.4, Theorem 3.23, pp.28--29; radial slicing and
  projection/Fubini applications, p.29. Almost-everywhere
  rectifiability: Theorem 3.24, p.29.
\end{itemize}

Normalization agrees with \texttt{simon-\allowbreak{}gmt}. Two evident dimensional
misprints occur in §3.4: the coarea Jacobian should vanish when the
derivative has rank less than \textbf{m}, not n; the area half of
Theorem 3.23 needs outer measure \textbf{H\^{}n}, not H\^{}m. The coarea
half uses the correct fiber and target dimensions. Our area statements
use Simon's 2018 Theorem 3.3 and Remark 3.4(2), whose dimensions are
correct.

\subsection{\texorpdfstring{\texttt{hajlasz-\allowbreak{}analysis}}{hajlasz-analysis}}\label{hajlasz-analysis}

Piotr Hajłasz, \emph{Analysis I}, University of Pittsburgh, undated
author-hosted lecture notes.
\href{https://www.pitt.edu/~hajlasz/Notatki/Analysis\%20I.pdf}{Author's
PDF}.

\begin{itemize}
\tightlist
\item
  Theorem 103 and the Cantor staircase discussion, printed/PDF p.76: a
  continuous increasing function can have derivative zero almost
  everywhere while its endpoint increment is positive.
\item
  Absolute continuity and the fundamental theorem: §8.1, Theorem 106,
  printed/PDF pp.77--78.
\item
  Convolution in L1: Theorem 87, pp.66--67; smooth approximate
  identities: Theorem 92, pp.70--71.
\end{itemize}

The Cantor paragraph prints \texttt{f\textquotesingle{}(0)\ =\ 0\ a.e.};
the variable should be x. The example in this library proves its needed
staircase properties directly. Derivative values on its exceptional null
set may be assigned arbitrarily for Lebesgue integration, but those
assignments do not make the staircase absolutely continuous or
Lipschitz.

\subsection{\texorpdfstring{\texttt{conrad-\allowbreak{}diff}}{conrad-diff}}\label{conrad-diff}

Keith Conrad, \emph{Differentiating Under the Integral Sign}, University
of Connecticut, undated author-hosted expository notes.
\href{https://kconrad.math.uconn.edu/blurbs/analysis/diffunderint.pdf}{Author's
PDF}.

\begin{itemize}
\tightlist
\item
  A sufficient differentiation theorem: §12, Theorem 12.3, printed/PDF
  p.16.
\item
  Local parameter bounds illustrating the theorem: Table 2, p.16, and
  discussion on p.17.
\end{itemize}

Theorem 12.3 treats continuous integrands and parameter derivatives,
dominated by parameter-independent integrable majorants near the
parameter in question. Its infinite-interval integrals are stated as
improper integrals. The library supplies Lebesgue proofs via dominated
convergence of difference quotients; it does not silently equate
conditional improper integration with Lebesgue integrability.

\subsection{Reading formulas across
versions}\label{reading-formulas-across-versions}

Use the source IDs above, not a theorem number alone: Simon's 2014 and
2018 drafts number area and coarea differently. The older
\href{https://web.stanford.edu/class/math285/ts-gmt.pdf}{2014 course
draft}, printed p.54, formula 3.4, also has the incorrect target
dimension H\^{}m; the following weighted formula 3.5 has H\^{}n.~The
library cites the corrected 2018 area statement.

The convention in this directory is scalar integration against positive
measures. A parametrization's area Jacobian is nonnegative and counts
repeated coverings. An oriented vector flux additionally needs a chosen
normal and a signed dot product; it is not determined by the scalar
Hausdorff measure alone.

No source access, symbolic calculation, numerical test, or review
recorded here is a claim of proof-assistant verification.

\end{document}
